% Generated methods addendum for the completed design analysis.
\subsection{Completed Pre-shock and Tariff-weighted Design}
All capability constructs are fixed before the tariff shock. Baseline ECI and COI are country averages over 2015--2017. Historical income group is the modal World Bank analytical classification over those years, with a 2017 tie-break rule. Low, middle, and high complexity regimes are fixed pre-shock ECI terciles. The causal sample excludes the United States and China because they are the direct policy parties rather than third-country comparison units.

Our primary treatment is a continuous tariff-weighted US-market diversion opportunity. For each third country, we calculate the 2015--2017 share of its US exports in Section 301-affected HS6 products and weight each product by China's corresponding pre-shock share of US imports. The fixed product profile is multiplied by the annual-average additional Section 301 duty. This construct captures the tariff-induced opportunity to displace Chinese suppliers, not a duty paid by the third country. A raw affected-basket measure and a China-import product-overlap measure are reported as supplementary channels. Because BACI is HS6 and the legal schedules are HTS8 or more detailed, all treatment measures are any-covered-line HS6 approximations; the replication package reports coverage, overlap, line-count, and source-hash audits.

GVC stability is measured primarily as the adverse deviation from a country's 2015--2017 forward-linkage mean, \(\min(FwdLinkage_{it}-\overline{FwdLinkage}_{i,pre},0)\), so values closer to zero indicate that a country avoided an erosion of its pre-shock GVC position. Signed annual linkage change and absolute deviation are reported as sensitivity outcomes. Export recovery is the logarithm of current exports relative to the fixed 2015--2017 export baseline; with country fixed effects, this is equivalent to modelling log exports with a country-specific baseline removed. Partner diversification excludes the United States and China and is measured primarily as one minus the destination-share HHI, with entropy and effective-destination-count sensitivity outcomes.

For outcome \(Y_{it}\), the primary model is
\begin{equation}
Y_{it}=\beta_1 T_{it}+\beta_2(ECI_i^{pre}\times Post_t)+\beta_3(ECI_i^{pre}\times T_{it})+\alpha_i+\lambda_t+\varepsilon_{it},
\end{equation}
where \(T_{it}\) is the scaled tariff-weighted diversion treatment. Country and year fixed effects are included. Inference uses country-clustered standard errors and 999-replication Rademacher wild-cluster bootstrap p-values in the final run. A sensitivity permits the frozen pre-shock GDP per capita, trade openness, institutional-quality, and resource-rent profiles to have separate year effects, avoiding controls that may themselves be affected by the shock.

H6 is evaluated in a pooled segmented model at the pre-specified median of frozen ECI, with the required tests \(\beta_L<0\) and \(\beta_H-\beta_L>0\). An outcome-selected threshold grid and a 999-replication selection distribution are reported only as exploratory localization, not as primary confirmatory evidence. E1 and E2 use single pooled group-interaction models with formal coefficient-difference contrasts. E3 is estimated first using direct country GPR observations; a same-historical-income-prioritised KNN profile match based only on frozen pre-shock country characteristics provides a supplementary missing-data sensitivity. Benjamini--Hochberg q-values are calculated over the complete reported confirmatory, sensitivity, H6, E1--E3, and mechanism test family.
